\section*{Bab \ref{ch:b3:cell-cycle-apoptosis} --- Kendali Daur Sel dan Kematian Sel Terprogram}

\begin{solution}{exo:b3:cell-cycle-apoptosis:1}
G1, S, G2, M. Kemajuan G1 dan \omterm{def:b3:cell-cycle-apoptosis:checkpoints}{titik batasnya}: \omterm{def:b3:cell-cycle-apoptosis:cdk}{siklin} D--Cdk4/6;
G1/S: \omterm{def:b3:cell-cycle-apoptosis:cdk}{siklin} E--Cdk2; S: \omterm{def:b3:cell-cycle-apoptosis:cdk}{siklin} A--Cdk2; G2/M: \omterm{def:b3:cell-cycle-apoptosis:cdk}{siklin} B--Cdk1.
\omterm{def:b3:cell-cycle-apoptosis:cdk}{Kompleks pemacu anafase} (APC/C) memusnahkan \omterm{def:b3:cell-cycle-apoptosis:cdk}{siklin} B dan sekurin lalu
mengakhiri mitosisnya.
\end{solution}

\begin{solution}{exo:b3:cell-cycle-apoptosis:2}
Hartwell: mutan \emph{cdc} khamir bertunas, yang terhenti pada tahap
tertentu, serta gagasan Start. Nurse: \emph{cdc2} khamir belah
sebagai kinase yang menakar waktu mitosisnya, dan homolog manusianya (Cdk1)
yang menyelamatkan khamirnya --- yakni kelestariannya. Hunt: \omterm{def:b3:cell-cycle-apoptosis:cdk}{siklin} pada telur
bulu babi, yakni protein yang disintesis sepanjang daurnya lalu dimusnahkan
pada tiap pembelahan.
\end{solution}

\begin{solution}{exo:b3:cell-cycle-apoptosis:3}
\omterm{def:b3:cell-cycle-apoptosis:checkpoints}{Titik batas} (G1 akhir): faktor pertumbuhan, hara, ukuran; bekerja pada
\omterm{def:b3:cell-cycle-apoptosis:cdk}{siklin} D--Cdk4/6 dan Rb. G2/M: DNA yang utuh dan telah bereplikasi; \omterm{def:b3:dna-repair:ddr}{ATM}/ATR
menghambat Cdc25 sehingga menjaga Cdk1 terfosforilasi dan tak aktif. Perakitan
gelendong: tiap kinetokornya melekat; kinetokor yang tak melekat membuat
kompleks Mad2--Cdc20 yang menghambat APC/C.
\end{solution}

\begin{solution}{exo:b3:cell-cycle-apoptosis:4}
\omterm{def:b3:cell-cycle-apoptosis:apoptosis}{Apoptosis}: selnya mengerut, \omterm{def:b3:chromatin-epigenetics:nucleosome}{kromatinnya} memadat, DNA-nya dipotong menjadi
sebuah tangga, membrannya menggelembung tetapi tetap tersegel,
\omterm{def:b3:cell-cycle-apoptosis:apoptosis}{fosfatidilserinanya} terpapar, mayatnya ditelan, dan tanpa peradangan.
\omterm{def:b3:cell-cycle-apoptosis:apoptosis}{Nekrosis}: selnya membengkak, membrannya pecah, isinya merembes, DNA-nya
terurai secara acak, dan timbul peradangan. Algojonya: \omterm{def:b3:cell-cycle-apoptosis:apoptosis}{kaspase}, yakni protease sisteina yang
memotong sesudah aspartat.
\end{solution}

\begin{solution}{exo:b3:cell-cycle-apoptosis:5}
Naik dari $5$ ke $30$ pada $1$ tiap menit: \qty{25}{min}. Penguraian dari
$30$ ke $5$ dengan waktu paruh \qty{2}{min}: $\log_{2}(30/5)\times 2 =
\qty{5.2}{min}$. Periodenya sekitar \qty{30}{min}, mitosisnya \qty{17}{\%} dari itu.
\end{solution}

\begin{solution}{exo:b3:cell-cycle-apoptosis:6}
(a) $C$ tak pernah dapat turun di bawah $C_{1}$: terkunci di dalam mitosis pada
cabang tingginya. (b) Tanpa fosforilasi penghambat: ambang $C_{2}$-nya
turun, sehingga mitosisnya mulai lebih awal pada ukuran kecil (fenotipe
``wee''). (c) Cdk1 tetap terfosforilasi: cabang tingginya tak terjangkau,
sehingga terhenti pada G2 dengan sel yang memanjang. (d) Seperti (b), dan
\omterm{def:b3:cell-cycle-apoptosis:checkpoints}{titik periksa} G2/M, yang bekerja lewat Wee1/Cdc25, tak lagi dapat menahan selnya.
\end{solution}

\begin{solution}{exo:b3:cell-cycle-apoptosis:7}
Nukleus G1 di dalam sitoplasma berfase S: titik asalnya telah diizinkan dan
sitoplasmanya memasok \omterm{def:b3:cell-cycle-apoptosis:cdk}{siklin} E/A--Cdk2 yang aktif, sehingga ia langsung masuk
S. Nukleus G2 di dalam sitoplasma berfase S: titik asalnya telah menyala dan
tak diizinkan ulang (faktor pengizinnya dimusnahkan atau dihambat CDK),
sehingga ia tak dapat bereplikasi lagi; ia menunggu sampai pasangannya mencapai
G2 lalu keduanya masuk mitosis bersama.
\end{solution}

\begin{solution}{exo:b3:cell-cycle-apoptosis:8}
Kinetokor yang tak melekat itu sebuah katalis: ia terus-menerus mengubah Mad2
menjadi konformasi aktifnya, sehingga menghasilkan penghambat Cdc20 yang dapat
berdifusi yang menyebar ke seluruh selnya lalu menjaga APC/C tetap mati di
mana-mana --- satu kinetokor saja sudah cukup. Tanpa Mad2, APC/C diaktifkan
begitu Cdk1 menyalakannya, entah kromosomnya telah melekat entah belum:
anafasenya mulai dengan kromosom yang tak melekat, sel anaknya menjadi
aneuploid, dan organismenya (atau galur selnya) mati kehilangan kromosom.
\end{solution}

\begin{solution}{exo:b3:cell-cycle-apoptosis:9}
Tanpa \emph{ced-9}: sel yang semestinya hidup justru mati --- kematian embrio
karena terlalu banyak kematian. Tanpa \emph{ced-3}: tak satu pun dari $131$
kematiannya terjadi, selnya bertahan lalu berdiferensiasi, dan cacingnya hidup.
Tanpa keduanya: tanpa kematian --- yakni fenotipe \emph{ced-3}. Karena
membuang algojonya meniadakan pengaruh membuang penjaganya, CED-9 bekerja di
hulu, dengan menahan CED-4/CED-3 tetap tak aktif; ketika CED-9 hilang keduanya
aktif, kecuali bila keduanya pun tak ada.
\end{solution}

\begin{solution}{exo:b3:cell-cycle-apoptosis:10}
Dengan $A = aC$, $\mathrm{d}C/\mathrm{d}t = k_{s} - k_{d}aC^{2}$, yakni
persamaan satu variabel dengan keadaan tunak $C^{*} = \sqrt{k_{s}/(k_{d}a)}$
dan kemiringan $-2k_{d}aC^{*} < 0$ di situ: simpangan apa pun meluruh
secara monoton, dan satu variabel orde satu tak dapat berosilasi. Kurva
berbentuk S memberi dua cabang mantap yang dipisahkan sebuah zona terlarang,
sehingga keadaannya harus melompat dan tak dapat menetap; sebuah tundaan yang
eksplisit pada umpan baliknya akan memberi osilasi lewat rute yang berbeda ---
remnya tiba terlambat, melampaui batas, dan seterusnya --- seperti pada banyak
irama hormon.
\end{solution}

\begin{solution}{exo:b3:cell-cycle-apoptosis:11}
Tiap pengindera yang tiba ditangkap seorang penjaga, sehingga tak terjadi
apa-apa sampai seluruh $N$ penjaganya terisi, yakni pada $t \approx N/r$;
pengindera berikutnya menjadi bebas, mengaktifkan Bax/Bak, yang pembentukan
porinya naik tajam terhadap cacahnya, lalu mitokondrianya membuka dalam
hitungan menit. Rangsangan yang lebih kuat hanya memendekkan waktunya lewat
$r$; di atas ambangnya hasilnya sama. Venetoklaks menduduki penjaganya
sehingga mengurangi $N$ efektifnya: waktunya menciut, dan sel yang tersiap ---
yang sudah dekat $N$ --- mati seketika.
\end{solution}

\begin{solution}{exo:b3:cell-cycle-apoptosis:12}
Bcl-2 yang berlebih menghalangi \omterm{def:b3:cell-cycle-apoptosis:pathways}{jalur intrinsiknya}: sel~B yang semestinya mati
ketika isyarat pertumbuhannya berakhir justru bertahan, dan populasinya tumbuh
karena kegagalan mati, pada laju lambat saat sel semacam itu dibuat ---
lamban, sampai sebuah cacat kedua menambahkan proliferasi. Hilangnya Rb
membuang \omterm{def:b3:cell-cycle-apoptosis:checkpoints}{titik batasnya}: pendahulu retinanya melanjut melewati daurnya
tanpa faktor pertumbuhan lalu membelah secepat yang diizinkan mesinnya ---
ganas. Kedua tumor itu adalah kedua program yang masing-masing gagal
sendirian: sebuah kendali kematian dan sebuah kendali pembelahan, dan
hilangnya salah satu saja sudah memadai bagi sebuah tumor, dengan yang kedua
lebih cepat.
\end{solution}

\begin{solution}{pb:b3:cell-cycle-apoptosis:1}
\textbf{1.} $10^{7}\times 3500/5 = 7\times 10^{9}$ sel tiap hari, kira-kira
\num{80000} tiap detik.
\textbf{2.} $3500/5 = 700$ tiap kripta tiap hari; $22\times 2^{5} = 704$.
\textbf{3.} $80\times 365 \approx \num{29000}$ pembelahan; $\times 0.6
\approx \num{17500}$ mutasi.
\textbf{4.} Sel transitnya dilepaskan dalam hitungan hari sambil membawa
mutasinya pergi; sel puncanya bertahan seumur hidup, sehingga tiap mutasinya
tersimpan dan tiap pembelahannya menambahkan lagi --- berdaur lambat
memperkecilnya.
\textbf{5.} Vilinya kosong dalam kira-kira $5$ hari masa transit normalnya;
membangunnya kembali memakan sehari pemulihan, lima pembelahan tiap
\qty{12}{h}, dan perpindahan naik ke vilinya --- empat sampai lima hari, dan
selama itu sawarnya gundul: pembelahan transit yang cepat pada ususnya
menjadikannya korban awal radiasi.
\textbf{6.} Mayat apoptotik yang tersegel tidak melepaskan isinya ke dalam
lumen yang penuh bakteri dan tidak menimbulkan peradangan; \omterm{def:b3:cell-cycle-apoptosis:apoptosis}{nekrosis} akan
menumpahkan enzim dan isyarat lalu meradangkan mukosanya pada tiap sel yang
dilepaskan.
\textbf{7.} $(40 - 8)/2 = \qty{16}{h}$.
\textbf{8.} $\log_{2}(40/8)\times \qty{10}{min} = 2.3\times 10 =
\qty{23}{min}$.
\textbf{9.} Periodenya sekitar \qty{16.4}{h}; Cdk1 aktif \qty{2.3}{\%}
waktunya.
\textbf{10.} Bukan $k_{s}$ saja melainkan lamanya ditahan pada ambangnya:
sel somatiknya menunggu pada \omterm{def:b3:cell-cycle-apoptosis:checkpoints}{titik batasnya} untuk faktor pertumbuhan dan pada
G2/M untuk DNA-nya, dan sel puncanya lebih lama daripada sel transitnya. Telur
kataknya tak memiliki \omterm{def:b3:cell-cycle-apoptosis:checkpoints}{titik periksa} dan sitoplasmanya berbekal segalanya,
sehingga hanya osilator telanjangnya yang menetapkan periodenya.
\textbf{11.} Pada cabang rendahnya dengan kadar \omterm{def:b3:cell-cycle-apoptosis:cdk}{siklin} di atas $C_{2}$ yang
normal: \omterm{def:b3:cell-cycle-apoptosis:checkpoints}{titik periksanya} telah menggeser kurva S-nya ke kanan dengan
menghambat Cdc25, sehingga lompatannya tak terjadi. Begitu diperbaiki, Cdc25
dilepaskan, ambangnya turun di bawah \omterm{def:b3:cell-cycle-apoptosis:cdk}{siklin} yang telah tertimbun, dan
selnya masuk mitosis seketika --- dan itulah sebabnya sel yang dilepaskan dari
\omterm{def:b3:cell-cycle-apoptosis:checkpoints}{titik periksa} masuk mitosis serentak.
\textbf{12.} \omterm{def:b3:cell-cycle-apoptosis:cdk}{Siklin} B maupun sekurinnya tidak diuraikan: selnya tetap berada
di metafase dengan kohesinnya utuh dan Cdk1 aktif, tanpa batas; ia mati
di situ atau akhirnya menyelinap keluar dengan genom yang tak terpilah.
\textbf{13.} $10^{11}$ pembelahan tiap hari.
\textbf{14.} Dengan \omterm{def:b3:cell-cycle-apoptosis:checkpoints}{titik periksanya} $10^{11}/5\times 10^{4} = 2\times 10^{6}$
sel anak aneuploid tiap hari; tanpanya, $2\times 10^{9}$.
\textbf{15.} $2\times 10^{4}$ tiap hari; sepanjang $80$ tahun $6\times 10^{8}$.
\textbf{16.} $10^{9}/50 = 2\times 10^{7}$ kesalahan pemilahan tiap hari:
setiap hari tumornya mencoba dua puluh juta kariotipe baru, dan seleksi
menyimpan yang tumbuh lebih cepat atau kebal terhadap sebuah obat.
\textbf{17.} Dengan setiap pembelahan salah memilah, tak ada garis keturunan
sel embrio yang mempertahankan genom euploid dan embrionya mati. Hilangnya
sebagian memberi aneuploidi sesekali pada sel yang bertahan (terutama tanpa
p53), yakni ketakmantapan kromosom yang memasok keragaman bagi sebuah tumor
tanpa membunuh organismenya.
\textbf{18.} Sebuah galat meiosis menaruh kromosom tambahannya ke dalam
gametnya, sehingga zigotnya dan setiap sel turunannya menjadi trisomik;
sebuah galat mitosis terjadi pada satu sel somatik lalu hanya diwarisi klonnya.
\textbf{19.} $10^{11}\times \qty{1}{ng} = \qty{100}{g}$ tiap hari,
\qty{36}{kg} tiap tahun --- kira-kira separuh massa tubuhnya tiap tahun.
\textbf{20.} \qty{20}{g} protein tiap hari, kira-kira sepertiga asupan
makanannya dan beberapa persen pergantian protein total tubuhnya.
\textbf{21.} $10^{11}$ mayat pada $24$ tiap makrofag tiap hari: $4\times
10^{9}$ makrofag penuh waktu, yakni \qty{4}{\%} waktu $10^{11}$
makrofagnya.
\textbf{22.} Mayat yang melisis melepaskan protease, DNA, ATP, dan isyarat
bahaya lain yang menimbulkan peradangan dan dapat memicu autoimunitas
terhadap antigen inti; mayat yang utuh mengiklankan dirinya dengan
\omterm{def:b3:cell-cycle-apoptosis:apoptosis}{fosfatidilserina} pada lembaran luarnya (dan dengan kehilangan isyarat ``jangan
makan aku''), lalu menarik fagosit dengan nukleotida yang dilepaskannya.
\textbf{23.} \omterm{def:b3:cell-cycle-apoptosis:pathways}{Jalur intrinsiknya} terhalang pada mitokondrianya. Rute
ekstrinsiknya masih bekerja di tempat kaspase-8 mengaktifkan kaspase-3
langsung, walaupun penguatannya lewat Bid hilang. Selnya
menumpuk karena tidak mati, bukan karena membelah lebih cepat --- tumbuh
lambat --- sampai sebuah cacat kemudian (kerap \emph{MYC}) menambahkan proliferasi.
\textbf{24.} \omterm{def:b3:cell-cycle-apoptosis:apoptosis}{Apoptosis} memakan beberapa jam, menuntut ATP, dan melewati
langkah yang dapat dihalangi --- penghambat \omterm{def:b3:cell-cycle-apoptosis:apoptosis}{kaspase}, ambang mitokondrianya ---
sehingga neuron di daerah remangnya dapat diselamatkan bila diobati tepat
waktu; neuron di intinya mati karena kegagalan energi dalam hitungan menit,
tanpa program yang dapat disela.
\textbf{25.} Periode osilator telanjangnya sekitar \qty{16}{h}; sekitar $2\times
10^{4}$ sel aneuploid yang membelah tiap hari di dalam tubuh yang normal;
\qty{100}{g} sel yang didaur ulang \omterm{def:b3:cell-cycle-apoptosis:apoptosis}{apoptosis} tiap hari.
\end{solution}
